\documentclass[11pt,letterpaper]{article}
\usepackage{../../Common/course_notes}
\usepackage{../../Common/course_macros}

\begin{document}
\setlength{\abovedisplayskip}{6pt}
\setlength{\belowdisplayskip}{6pt}
\setlength{\textfloatsep}{10pt}
\setlength{\intextsep}{8pt}
\color{blackmain}

\assignmentheader{MÓDULO 00}{SEMANA 03}{Tarea 3: El Grupo de Rotaciones $SO(2)$, $SO(3)$ y Traslaciones}

\noindent\textbf{Instrucciones:} Desarrolla de manera clara y justificada cada uno de los problemas siguientes. Apóyate en esquemas geométricos cuando sea apropiado y fundamenta tus deducciones en las definiciones formales y leyes de transformación vistas.

% =============================================================================
\manualsection{Parte A: Axiomas de Grupo, Ortogonalidad y Reflexiones}

\begin{infobox}{Problema 1: Verificación de los Axiomas de Grupo para $SO(2)$}
Considera el conjunto de matrices de rotación en el plano:
\[
SO(2) = \left\{ R(\theta) = \begin{pmatrix} \cos\theta & -\sin\theta \\ \sin\theta & \cos\theta \end{pmatrix} \;\middle|\; \theta \in [0, 2\pi) \right\}.
\]
\begin{enumerate}[label=(\alph*)]
    \item Demuestra la \textbf{cerradura}: calcula el producto matricial $R(\theta_1)R(\theta_2)$ y comprueba que coincide exactamente con $R(\theta_1 + \theta_2)$ usando identidades trigonométricas de suma de ángulos.
    \item Identifica el elemento neutro $I$ y determina para qué valor del ángulo $\theta$ se obtiene.
    \item Para cualquier $R(\theta)$, calcula su matriz transpuesta $R(\theta)^T$. Verifica explícitamente que $R(\theta)^T R(\theta) = I$, y deduce que el inverso es $R(\theta)^{-1} = R(-\theta) \in SO(2)$.
    \item Demuestra que $R(\theta_1)R(\theta_2) = R(\theta_2)R(\theta_1)$, concluyendo que $(SO(2), \cdot)$ es un grupo abeliano (conmutativo).
\end{enumerate}
\end{infobox}

\begin{infobox}{Problema 2: Rotaciones Propias frente a Reflexiones e Inversiones}
Considera las matrices en $\R^3$:
\[
R_1 = \begin{pmatrix} 0 & -1 & 0 \\ 1 & 0 & 0 \\ 0 & 0 & 1 \end{pmatrix}, \qquad
S = \begin{pmatrix} 1 & 0 & 0 \\ 0 & 1 & 0 \\ 0 & 0 & -1 \end{pmatrix}, \qquad
P = -I = \begin{pmatrix} -1 & 0 & 0 \\ 0 & -1 & 0 \\ 0 & 0 & -1 \end{pmatrix}.
\]
\begin{enumerate}[label=(\alph*)]
    \item Verifica si $R_1$, $S$ y $P$ pertenecen al grupo ortogonal $O(3)$ evaluando sus productos con la transpuesta.
    \item Calcula los determinantes $\det(R_1)$, $\det(S)$ y $\det(P)$. ¿Cuáles pertenecen a $SO(3)$?
    \item Explica por qué la inversión espacial completa $P = -I$ en $\R^3$ invierte la orientación ($\det P = -1$), mientras que en $\R^2$ la matriz $-I_{2\times 2}$ tiene $\det(-I) = +1$ y es una rotación propia de $180^\circ$.
\end{enumerate}
\end{infobox}

% =============================================================================
\manualsection{Parte B: No-Conmutatividad, Fórmula de Rodrigues y Eje Invariante}

\begin{infobox}{Problema 3: No-Conmutatividad en Tres Dimensiones y Conmutadores}
Considera las rotaciones de $\pi/2$ alrededor de los ejes coordenados $X$ y $Z$:
\[
R_x(\pi/2) = \begin{pmatrix} 1 & 0 & 0 \\ 0 & 0 & -1 \\ 0 & 1 & 0 \end{pmatrix}, \qquad
R_z(\pi/2) = \begin{pmatrix} 0 & -1 & 0 \\ 1 & 0 & 0 \\ 0 & 0 & 1 \end{pmatrix}.
\]
\begin{enumerate}[label=(\alph*)]
    \item Calcula los productos matriciales $R_A = R_x(\pi/2)R_z(\pi/2)$ y $R_B = R_z(\pi/2)R_x(\pi/2)$.
    \item Aplica $R_A$ y $R_B$ al vector unitario $\hat{e}_3 = (0,0,1)^T$. Compara las posiciones finales y verifica que $R_A \neq R_B$.
    \item Calcula el conmutador matricial $[R_x(\pi/2), R_z(\pi/2)] \equiv R_x R_z - R_z R_x$.
\end{enumerate}
\end{infobox}

\begin{infobox}{Problema 4: Construcción de la Fórmula de Rotación de Rodrigues}
Sea $\hat{n} = \frac{1}{\sqrt{3}}(1, 1, 1)^T$ la dirección de un eje de giro y $\theta = \frac{2\pi}{3}$ ($120^\circ$) el ángulo de rotación.
\begin{enumerate}[label=(\alph*)]
    \item Construye la matriz antisimétrica $K$ asociada a $\hat{n}$, dada por $K_{ij} = -\epsilon_{ijk} n_k$.
    \item Calcula $K^2$ y comprueba que para cualquier vector $\vec{v}$ se cumple $K^2 \vec{v} = \hat{n}(\hat{n}\cdot\vec{v}) - \vec{v}$.
    \item Utiliza la fórmula matricial de Rodrigues $R(\hat{n}, \theta) = I + (\sin\theta)K + (1 - \cos\theta)K^2$ para calcular explícitamente la matriz de rotación $R(\hat{n}, \theta)$.
    \item Aplica la matriz obtenida al vector $(1, 0, 0)^T$ y describe geométricamente la permutación cíclica de los ejes coordenados.
\end{enumerate}
\end{infobox}

% =============================================================================
\manualsection{Parte C: Transformaciones Activas frente a Pasivas y Dinámica}

\begin{infobox}{Problema 5: Rotaciones Activas frente a Pasivas}
En un plano con base canónica $\{\hat{e}_1, \hat{e}_2\}$, se considera el vector físico $\vec{v} = 2\hat{e}_1 + \hat{e}_2$.
\begin{enumerate}[label=(\alph*)]
    \item \textbf{Rotación Activa:} El vector se rota un ángulo $\theta = 45^\circ$ en sentido antihorario manteniendo los ejes fijos. Calcula las componentes del nuevo vector rotado $\vec{v}' = R(45^\circ)\vec{v}$.
    \item \textbf{Rotación Pasiva:} El vector físico $\vec{v}$ permanece fijo en el espacio, pero los ejes coordenados se rotan un ángulo $\theta = 45^\circ$ en sentido antihorario: $\hat{e}'_i = R(45^\circ)\hat{e}_i$.
    Calcula las nuevas componentes $[v']$ del vector en la base rotada usando $[v'] = R(45^\circ)^{-1}[v] = R(45^\circ)^T [v]$.
    \item Dibuja ambos casos en diagramas cartesianos separados y explica por qué las componentes numéricas en una rotación pasiva de ángulo $\theta$ coinciden con las de una rotación activa de ángulo $-\theta$.
\end{enumerate}
\end{infobox}

\begin{infobox}{Problema 6: Rotaciones Infinitesimales y el Vector Velocidad Angular}
\begin{enumerate}[label=(\alph*)]
    \item Considera una rotación dependiente del tiempo $R(t)$ que parte del reposo con $R(0) = I$. A partir de la identidad de rigidez $R(t)^T R(t) = I$, deriva respecto al tiempo en $t=0$ y demuestra que la matriz $\Omega \equiv \dot{R}(0)$ es estrictamente antisimétrica ($\Omega^T = -\Omega$).
    \item Si $\Omega = \begin{pmatrix} 0 & -\omega_3 & \omega_2 \\ \omega_3 & 0 & -\omega_1 \\ -\omega_2 & \omega_1 & 0 \end{pmatrix}$ y $\vec{\omega} = (\omega_1, \omega_2, \omega_3)^T$, demuestra mediante multiplicación directa que para cualquier posición $\vec{r} = (x, y, z)^T$:
    \[
    \Omega\vec{r} = \vec{\omega} \times \vec{r}.
    \]
    \item Demuestra que para cualquier vector $\vec{x}$, el producto escalar $\vec{x} \cdot (\Omega\vec{x}) = 0$. Explica por qué este resultado garantiza físicamente que la distancia de una partícula al origen permanece constante en una rotación pura.
\end{enumerate}
\end{infobox}

\begin{infobox}{Problema 7: El Grupo Euclídeo Especial $SE(2)$ y Coordenadas Homogéneas}
En robótica y mecánica analítica, un movimiento rígido en el plano combina una rotación de ángulo $\theta$ con una traslación $\vec{a} = (a_1, a_2)^T$. Representamos esta operación afín como una matriz de $3 \times 3$ actuando sobre coordenadas homogéneas:
\[
M(\theta, \vec{a}) = \begin{pmatrix} \cos\theta & -\sin\theta & a_1 \\ \sin\theta & \cos\theta & a_2 \\ 0 & 0 & 1 \end{pmatrix}.
\]
\begin{enumerate}[label=(\alph*)]
    \item Multiplica dos matrices de este tipo, $M(\theta_2, \vec{a}_2) M(\theta_1, \vec{a}_1)$, y demuestra que el producto reproduce la regla de composición de $SE(2)$:
    \[
    R_{\text{total}} = R(\theta_1 + \theta_2), \qquad \vec{a}_{\text{total}} = \vec{a}_2 + R(\theta_2)\vec{a}_1.
    \]
    \item Demuestra que la matriz inversa en coordenadas homogéneas está dada por:
    \[
    M(\theta, \vec{a})^{-1} = \begin{pmatrix} R(\theta)^T & -R(\theta)^T \vec{a} \\ \vec{0}^T & 1 \end{pmatrix}.
    \]
    \item Aplica la matriz $M(\theta, \vec{a})$ a un punto $P = (x, y, 1)^T$ y a un vector de desplazamiento $\vec{v} = (v_x, v_y, 0)^T$. Explica a partir de la cuarta fila por qué la traslación $\vec{a}$ desplaza la ubicación del punto pero deja inalterado al vector relativo.
\end{enumerate}
\end{infobox}

\end{document}
